\section{METHOD AND MATERIALS}\label{sec:method}
\paragraph{Processing pipeline overview.}
\label{sec:pipeline_overview}

The proposed pipeline is an epoch-aware, threshold-based pipeline that transforms a continuous multichannel EEG recording into a non-overlapping epoch through a
fixed sequence of five operations.
The Figure~\ref{fig:pipeline_flow} summarises this
sequence as a flow chart, from the preprocessing, epoching, and the double thresholding which is the proposed method, through three detection stages (A, B, and C) which we will discussed in detail in section.




\begin{figure*}[!t]
    \centering
    \includegraphics[width=\textwidth]{flowchart}
    \caption{Epoch-aware double-thresholding pipeline for multichannel EEG blink-region detection. Continuous EEG signals from FP1, AF3, F3, FP2, AF4, and F4 are downsampled to 100 Hz, band-pass filtered from 1–20 Hz, and segmented into non-overlapping 30-s epochs after low-health epochs are removed. Stage A applies channel-specific peak-to-peak screening and a multichannel union rule to select suspicious epochs; E2, E5, and E7 form the candidate set in this example. Stage B pools and concatenates samples from these candidate epochs to estimate a robust sample-level threshold separately for each channel. Stage C applies the resulting channel-specific threshold to all valid epochs, including candidate and non-candidate epochs, and reports each detected blink region by epoch ID, channel, onset time, and duration. Orange denotes Stage A candidate epochs, while red denotes detected blink regions.}
    \label{fig:pipeline_flow}
\end{figure*}


\paragraph{Preprocessing}
First, the continuous recording is downsampled to a common sampling grid of \(100\)~Hz. Standardising the sampling rate makes the
amplitude thresholds and the minimum-duration gate comparable across recordings acquired on different hardware, and reduces the per-sample computational cost of the downstream stages without discarding any information needed to resolve a blink, whose
spectral content lies well below the new Nyquist frequency.

Second, a zero-phase band-pass filter from \(1.0\)~Hz  to \(20.0\)~Hz
is applied which is removes slow drifts and DC offsets at the low end and high-frequency muscle and line noise at the high end, isolating the low-frequency, high-amplitude deflection that characterises an eye blink.

Third, each filtered recording is split into non-overlapping, fixed-length epochs. Epoching is what makes the thresholds
\emph{adaptive}, rather than deriving a single level for the entire session, the pipeline treats each short epoch as a local statistical context.


%Finally, the detected regions are compared with the
%ground-truth annotations at the \emph{event} level using overlap-based matching, yielding
%precision, recall, and \(F_1\); these are macro-averaged across sessions and compared between
%configurations with Wilcoxon signed-rank tests, as detailed in
%Section~\ref{sec:eval}.

%Recordings were segmented into fixed-length epochs, with 30~s used as the primary epoch length and 10, 20, 40, and 60~s additionally evaluated in an epoch-sensitivity sweep. The proposed detector was implemented as a three-stage epoch-aware threshold pipeline. First, suspicious/bad epochs were screened using an autoreject-based procedure with a fixed random state and a requirement that at least one epoch be flagged. Second, a sample-level amplitude threshold was estimated from the screened epochs using a robust centre and a standard-deviation multiplier of 3.5. Third, threshold crossings were expanded into blink regions. The detailed equations for each stage were given in the following subsections rather than duplicated in this overview.

\subsection{Selection of Bad Epochs}

The first stage of the proposed pipeline performs epoch-level screening using an
autoreject-inspired threshold selection procedure. Its purpose is not to estimate
the onset or offset of a blink, but to identify epochs whose amplitudes are
sufficiently abnormal to warrant subsequent sample-level localization. Let
\(X \in \mathbb{R}^{N \times C \times T}\) denote the epoched multichannel
signal after preprocessing, where \(N\) is the number of valid epochs, \(C\) is
the number of channels, and \(T\) is the number of samples per epoch. The sample
at epoch \(i \in \mathcal{E}=\{1,\ldots,N\}\), channel
\(c \in \mathcal{C}=\{1,\ldots,C\}\), and within-epoch time index
\(t \in \{1,\ldots,T\}\) is denoted by \(x_{i,c,t}\).

For each epoch and channel, an amplitude-range statistic is computed using the
peak-to-peak value
\[
a_{i,c}
= \max_{1 \leq t \leq T} x_{i,c,t}
- \min_{1 \leq t \leq T} x_{i,c,t}.
\]
Large values of \(a_{i,c}\) indicate that the corresponding epoch contains a
high-amplitude excursion in channel \(c\), a pattern consistent with ocular or
other transient artifacts. Rather than imposing a fixed rejection level, the
pipeline estimates a channel-specific threshold \(\tau_c^\star\) from the data.
Candidate thresholds are taken from the empirical range of peak-to-peak
amplitudes for that channel, denoted
\(\mathcal{Q}_c=\{a_{i,c}: i \in \mathcal{E}\}\), and are evaluated by a
cross-validation criterion.

Specifically, let \(\{(\mathcal{T}_k,\mathcal{V}_k)\}_{k=1}^{K}\) be
cross-validation splits of the epoch indices into training sets
\(\mathcal{T}_k\) and validation sets \(\mathcal{V}_k\). For a candidate
threshold \(\tau \in \mathcal{Q}_c\), the training epochs retained for channel
\(c\) in fold \(k\) are
\[
\mathcal{G}_{k,c}(\tau)
= \{ i \in \mathcal{T}_k : a_{i,c} \leq \tau \}.
\]
These retained epochs define a threshold-dependent training template
\[
\bar{x}^{(\tau)}_{k,c}(t)
= \frac{1}{|\mathcal{G}_{k,c}(\tau)|}
\sum_{i \in \mathcal{G}_{k,c}(\tau)} x_{i,c,t},
\]
where candidates yielding an empty retained set are treated as inadmissible. The
corresponding validation reference is computed as the pointwise median across
the validation epochs,
\[
\tilde{x}_{k,c}(t)
= \operatorname{median}_{i \in \mathcal{V}_k} x_{i,c,t}.
\]
The use of a median validation reference makes the criterion less sensitive to
large validation-set artifacts. The validation loss for channel \(c\) is then
defined as
\[
\mathcal{J}_{c}(\tau)
= \frac{1}{K}\sum_{k=1}^{K}
\left[
\frac{1}{T}\sum_{t=1}^{T}
\left\{\bar{x}^{(\tau)}_{k,c}(t)-\tilde{x}_{k,c}(t)\right\}^{2}
\right]^{1/2}.
\]
The selected channel threshold is the candidate that minimizes the average
validation discrepancy,
\[
\tau_c^\star
= \arg\min_{\tau \in \mathcal{Q}_c} \mathcal{J}_{c}(\tau).
\]
In practice, this objective may be searched by sequential Bayesian optimization
or by a random search over candidate peak-to-peak thresholds; both choices
retain the same validation objective and rejection rule.

After the thresholds \(\{\tau_c^\star\}_{c=1}^{C}\) have been estimated, each
epoch receives a channel-level artifact indicator
\[
z_{i,c}
= \mathbf{1}\{a_{i,c} > \tau_c^\star\},
\]
where \(\mathbf{1}\{\cdot\}\) is the indicator function. The present pipeline
uses the union of the channel-level decisions to form the epoch-level screening
mask,
\[
b_i
= \mathbf{1}\left\{\sum_{c=1}^{C} z_{i,c} \geq 1\right\}
= \mathbf{1}\left\{\exists c \in \mathcal{C}: a_{i,c} > \tau_c^\star\right\}.
\]
The set of suspicious epochs is therefore
\[
\mathcal{B}=\{i \in \mathcal{E}: b_i=1\}.
\]
These epochs are interpreted as likely blink-contaminated or artifact-heavy
segments. They provide the restricted data subset from which the next stage
estimates sample-level blink-region thresholds.

\paragraph{Computation of Sample-Level Blink-Region Threshold Candidates}

Stage~B estimates channel-specific sample-level amplitude thresholds using signal samples drawn from the candidate epochs \(\mathcal{B}\) identified in Stage~A.
The estimation combines a measure of central tendency with a robust measure of dispersion and a sensitivity parameter, allowing the threshold to adapt to the amplitude characteristics of each channel while limiting the influence of isolated extreme samples.

%In Stage~B, a robust sample-level amplitude threshold is estimated using only the samples
%pooled from the flagged epochs \[\mathcal{B}\], as discussed in previous subsection., as
%\(\theta = \xi + k\,(1.4826 \cdot \mathrm{MAD})\), where the centre \(\xi\) is the
%\emph{median} in the primary configuration (or the \emph{mean} in the contrast
%configuration), \(\mathrm{MAD}\) is the median absolute deviation, the factor
%\(1.4826\) rescales the MAD to a standard-deviation-consistent dispersion, and \(k\) is
%the sensitivity multiplier. Estimating the threshold from the flagged epochs concentrates
%the calculation on blink-enriched signal, while the median-plus-MAD construction keeps the
%estimate stable in the presence of isolated extreme samples. If too few epochs are flagged,
%the estimate falls back to the full valid-epoch set so that the threshold is never derived
%from an unreliably small sample.
%
%
%The second stage converts the epoch-level artifact decision into a
%sample-level thresholding problem. The autoreject-style screening in the
%previous stage identifies a set of suspicious epochs \(\mathcal{B}\), but its
%thresholds \(\tau_c^\star\) are defined on peak-to-peak amplitudes over whole
%epochs. Such thresholds are therefore appropriate for deciding whether an epoch
%is abnormal, but they are not directly interpretable as sample-wise blink
%crossing levels. Stage B estimates a separate channel-specific threshold in
%the amplitude domain of the time series itself.

Let \(\mathcal{S}\) denote the set of epochs used for estimating the
sample-level threshold. In the primary analysis,
\(\mathcal{S}=\mathcal{B}\), so that the estimate is driven by epochs already
identified as likely to contain blink-related activity.

%If too few suspicious
%epochs are available, the method falls back to the full valid epoch set
%\(\mathcal{E}\) to avoid estimating a threshold from an insufficient sample:
%\[
%\mathcal{S}
%=
%\begin{cases}
%\mathcal{B}, & |\mathcal{B}| \geq n_{\min},\\
%\mathcal{E}, & |\mathcal{B}| < n_{\min},
%\end{cases}
%\]
%where \(n_{\min}\) is the minimum number of screened epochs required for a
%stable estimate.

For each channel \(c\), all samples from the selected epochs are pooled into a
one-dimensional empirical sample set
\[
\mathcal{Y}_c
= \{x_{i,c,t}: i \in \mathcal{S},\; 1 \leq t \leq T\}.
\]
%This pooling step marks the methodological transition from epoch-level rejection to sample-level localization: the quantities used in Stage B are no
%longer peak-to-peak ranges \(a_{i,c}\), but the individual signal amplitudes
%\(x_{i,c,t}\) that will later be thresholded to form candidate blink regions.

The sample-level threshold is estimated using a robust location-plus-dispersion form. Let \(\xi_c\) be the channel-specific center of the sample distribution.
The default choice is the median,
\[
\xi_c = \operatorname{median}_{y \in \mathcal{Y}_c} y,
\]

It should be note that, the arithmetic mean may be substituted when a mean-centered threshold
is required for comparison.

The dispersion is estimated by the median absolute
deviation (MAD),
\[
\operatorname{MAD}_c
= \operatorname{median}_{y \in \mathcal{Y}_c}
\left|y-\operatorname{median}_{u \in \mathcal{Y}_c}u\right|,
\]
and is rescaled to a standard-deviation-compatible quantity by
\[
\sigma^{\mathrm{rob}}_c = 1.4826\,\operatorname{MAD}_c .
\]
The constant \(1.4826\) is the usual normal-consistency factor for the MAD.

For a sensitivity parameter \(k>0\), the resulting sample-level blink-region
threshold candidate for channel \(c\) is
\[
\theta_c = \xi_c + k\sigma^{\mathrm{rob}}_c .
\]
The set
\[
\Theta = \{\theta_c: c \in \mathcal{C}\}
\]
constitutes the channel-wise threshold candidates passed to the blink-region extraction stage.

%This construction deliberately separates the statistical role of the two thresholds in the pipeline. The quantities \(\tau_c^\star\) from Stage A are  epoch-rejection thresholds defined on peak-to-peak amplitudes,
%whereas
%\(\theta_c\) is a sample-level amplitude threshold defined on the time-series  samples themselves.
%Estimating \(\theta_c\) from \(\mathcal{S}\) concentrates  the calculation on signal segments enriched for blink-like excursions, while  the robust center and MAD-scaled dispersion reduce the influence of isolated  extreme peaks within those same segments.

\subsection{Extraction of Blink Regions}
In Stage~C, the threshold is scanned across the signal in
the manner of the Kleifges crossing detector: maximal contiguous runs of supra-threshold
samples define candidate blink regions, brief crossings shorter than a minimum-duration gate
are discarded, and the accepted regions are mapped from their concatenated coordinates back
to epoch-relative onset and offset times.

The final stage applies the sample-level thresholds
\(\Theta=\{\theta_c\}_{c=1}^{C}\) to obtain channel-wise blink regions. For
each channel \(c\), the epochs selected for localization are ordered as
\(\mathcal{S}=\{i_1,\ldots,i_R\}\), where \(R=|\mathcal{S}|\). Their samples
are concatenated into a one-dimensional signal
\[
u_c[n] = x_{i_r,c,t},
\qquad
n=(r-1)T+t,\quad
r \in \{1,\ldots,R\},\quad
t \in \{1,\ldots,T\}.
\]
Thus \(u_c \in \mathbb{R}^{RT}\) is the channel-specific signal on which blink
onsets and offsets are estimated.

%If the analysis is run in the fallback mode  described above, \(\mathcal{S}\) is the full valid epoch set; otherwise it is the set of epochs identified as suspicious in Stage A.

Candidate blink samples are defined by a threshold-crossing indicator
\[
h_c[n] = \mathbf{1}\{u_c[n] > \theta_c\},
\qquad n=1,\ldots,RT .
\]
After any required polarity normalization, \(h_c[n]=1\) indicates that the
sample amplitude exceeds the channel-specific blink-region threshold. Blink
regions are then defined as maximal contiguous runs of supra-threshold samples.
A candidate interval for channel \(c\) is written as the half-open index range
\([s_{c,r},e_{c,r})\), where
\[
h_c[n]=1 \quad \text{for all } n=s_{c,r},\ldots,e_{c,r}-1,
\]
with either \(s_{c,r}=1\) or \(h_c[s_{c,r}-1]=0\), and either
\(e_{c,r}=RT+1\) or \(h_c[e_{c,r}]=0\). Equivalently, the onset
\(s_{c,r}\) is the first sample in a run that crosses above \(\theta_c\), and
the offset \(e_{c,r}\) is the first subsequent sample at which the signal has
returned to the sub-threshold state.

%To suppress brief threshold crossings that are unlikely to represent complete
%blinks, each candidate interval is required to exceed a minimum duration. Let
%\(f_s\) be the sampling frequency and let \(\ell_{\min}\) be the minimum blink
%duration in seconds. The corresponding sample length is
%\[
%L_{\min}=\ell_{\min} f_s .
%\]
%The accepted blink-region set for channel \(c\) is
%\[
%\mathcal{R}_c
%=
%\left\{
%(s_{c,r},e_{c,r}) :
%e_{c,r}-s_{c,r} > L_{\min}
%\right\}.
%\]
%This duration gate is applied independently for each channel. Consecutive
%supra-threshold samples are merged into the same region by construction, while
%two runs separated by at least one sub-threshold sample are treated as distinct
%candidate blink regions.

Finally, the accepted intervals are mapped from concatenated coordinates back to epoch-local coordinates using the known epoch boundaries.

%If an accepted
%onset \(s_{c,r}\) lies in the concatenated segment corresponding to epoch
%\(i_q\), its within-epoch onset sample is
%\[
%\hat{s}_{c,r}=s_{c,r}-(q-1)T,
%\]
%and its onset time is \((\hat{s}_{c,r}-1)/f_s\) seconds relative to the
%beginning of that epoch. The interval duration is
%\[
%\hat{d}_{c,r}=\frac{e_{c,r}-s_{c,r}}{f_s}.
%\]
The final output of this stage is therefore a channel-indexed collection of
blink regions, each represented by an epoch identifier, a channel label, an
onset index or time, and an offset or duration.

%The extraction remains
%channel-wise throughout: the threshold \(\theta_c\) estimated for channel \(c\)
%is applied only to that channel's signal, and the resulting regions
%\(\mathcal{R}_c\) are retained as channel-specific blink candidates.



\subsection{Epoch Segmentation and Detection Pipeline}






\subsection{Problem Definition}\label{sec:eval}

Let a blink annotation be an interval \(g_i = [s_i, e_i]\) and a detector output be an interval \(d_j = [\hat{s}_j, \hat{e}_j]\). A detected region is counted as a true positive if it overlaps at least one unmatched ground-truth region. Unmatched detections are false positives; unmatched ground-truth regions are false negatives. This event-level treatment follows the logic of interval-based detection rather than pointwise accuracy, which is appropriate because blink boundaries are uncertain, annotation granularity varies, and non-event samples dominate the timeline \citep{saito2015precision,salles2024softed}.

\subsection{Datasets}
\label{sec:datasets}
Two naturalistic driving-EEG corpora were evaluated: Raja and Cao2018. The Raja corpus comprised 46 complete sessions recorded with a 128-channel EGI system (ANT Neuro), referenced to CPz under a modified international 10-20 layout, during a simulated lane-keeping driving task; ground-truth blink annotations were derived from human EAR/EOG labelling. The Cao2018 corpus was the sustained-attention driving corpus of \citet{cao2019multichannel}, in which 32-channel EEG was recorded during an approximately 90-minute lane-keeping task; 58 complete sessions were used, and 4 non-complete sessions were excluded. Because the published preprocessed version of Cao2018 removed eye blinks, the raw recordings were used and paired with external blink-region annotations.


\subsection{Preprocessing}
All detectors used identical preprocessing, consisting of a 1--20 Hz band-pass filter, no resampling, and EEG-channel selection. The channel scope was dataset specific: for Raja, the analysis was restricted to a minimal frontal montage comprising three channels (E3, E9, and E22), whereas for Cao2018, all 32 EEG channels were provided to the pipeline and the best channel was selected separately for each session. Thus, Raja channel selection was confined to frontal sites by construction, while Cao2018 channel selection was performed over the full montage.




\subsection{Detector Configurations}
Four detector configurations were compared under identical preprocessing and epoching. The proposed methods were Proposed-Med, which implemented the three-stage pipeline with a median centre, and Proposed-Mean, which used a mean centre. The conventional baselines were BLINKER-concat and MNE-annot.


\subsection{Experiment Overview}
\label{sec:experiment_overview}

The experiments follow a deliberate logical progression. Experiment~1 first selects the
best Stage-A channel-selection configuration; Experiment~2 then benchmarks that chosen
configuration against competing detectors; and Experiments~3--8 stress-test and
characterise the resulting pipeline along orthogonal axes (epoch duration, boundary
tolerance, fallback sensitivity, blink morphology, epoch-health exclusion, and long-blink
recall). Fixing the channel configuration first, before any baseline comparison, ensures
that the headline comparison in Experiment~2 uses a principled spatial selection rather
than an arbitrary default, and that every subsequent robustness check operates on the same
ideal configuration. Each experiment is summarised below in terms of its purpose,
motivation, and objective.

\paragraph{Experiment~1 --- Channel-selection strategies.}
\textit{Purpose:} determine which Stage-A channel-selection group yields the best detector,
evaluated on both datasets (Raja and Cao2018) and under both estimators (median and mean),
before any other experiment is run.
\textit{Motivation:} the current default rule flags an epoch whenever \emph{any} channel
over \emph{all} channels exceeds its peak-to-peak threshold, which can be triggered by a
single noisy non-ocular channel and inflate the number of flagged epochs; a more selective
spatial group may screen blinks more cleanly.
\textit{Objective:} run the complete Stage~A$\to$B$\to$C pipeline on each candidate
channel group and select the best group by weighing multiple criteria rather than a single
metric. The decision criteria are:
\begin{description}
  \item[(i) Downstream event $F_1$.] The end-to-end event-level $F_1$ of the full pipeline
    when restricted to the group; this is the ultimate quantity the detector is optimised
    for and is therefore the primary criterion.
  \item[(ii) Stage-A epoch-selection $F_1$.] The precision and recall of the
    suspicious-epoch selection itself, scored against epoch-level ground truth (an epoch is
    positive if it contains at least one annotated blink); this isolates the quality of the
    screening stage from the downstream localisation.
  \item[(iii) Number of channels for Stage~A.] The count of channels the group requires,
    as a measure of montage cost, simplicity, and portability; a group that performs
    comparably with fewer channels is preferred because it is cheaper to acquire and easier
    to deploy.
  \item[(iv) False-positive epoch rate / \% epochs flagged.] The fraction of non-blink
    epochs erroneously flagged and the overall proportion of epochs flagged; a low value
    indicates robustness to artefacts and avoids over-flagging that would dilute the
    Stage-B threshold estimate.
  \item[(v) Recall on safety-critical blinks.] The recall specifically on long
    drowsiness-related closures, reflecting a ``do-not-miss'' requirement: a configuration
    that maximises overall $F_1$ but systematically misses safety-critical blinks is
    undesirable.
  \item[(vi) Cross-dataset consistency.] The agreement of the chosen group's behaviour
    between Raja and Cao2018; a group that is best on one dataset but poor on the other does
    not generalise.
  \item[(vii) Best-channel stability.] The stability of which channel (or channels) drives
    detection across sessions and subjects; an unstable best channel signals an
    overfitted, fragile selection.
  \item[(viii) Estimator robustness.] The consistency of the group's ranking under the
    median versus mean centre at Stage~B; a group whose advantage survives both estimators
    is more trustworthy than one that depends on a particular centre.
\end{description}
The configuration that best balances these criteria is designated the ideal configuration
and is carried forward unchanged into Experiments~2--8.

\paragraph{Experiment~2 --- Strategy comparison.}
\textit{Purpose:} benchmark the proposed pipeline, configured with the ideal channel group
from Experiment~1, against competing detection algorithms.
\textit{Motivation:} a method is only convincing if it outperforms the established
alternatives under identical preprocessing, epoching, and scoring.
\textit{Objective:} compare BLINKER-concat, MNE-annot, DBO, Proposed-Mean, and Proposed-Med
at the event level on both datasets, reporting precision, recall, and $F_1$ with paired
significance tests, and isolating both the estimator contrast (median versus mean) and the
cross-dataset generalisation gap.

\paragraph{Experiment~3 --- Epoch-duration stability.}
\textit{Purpose:} characterise how sensitive the proposed pipeline is to the choice of
epoch length.
\textit{Motivation:} the \(30\)~s primary epoch is a design choice, and a robust method
should not depend critically on it.
\textit{Objective:} sweep the epoch duration (e.g.\ \(10/20/40/60\)~s around the \(30\)~s
reference) and quantify the resulting variation in event-level $F_1$, testing each duration
against the reference with Wilcoxon signed-rank tests.

\paragraph{Experiment~4 --- Boundary-tolerance (IoU) stability.}
\textit{Purpose:} assess how the results depend on how strictly a detection must overlap a
ground-truth interval to count as a match.
\textit{Motivation:} blink boundaries are inherently uncertain, so the matching tolerance
is a scoring choice that could materially change the reported numbers.
\textit{Objective:} recompute event-level $F_1$ across a range of intersection-over-union
matching thresholds (e.g.\ \(\{0.0, 0.1, 0.2, 0.3, 0.5\}\)) and report the range of
$F_1$ as a measure of boundary-tolerance sensitivity.

\paragraph{Experiment~5 --- $n_{\min}$ / Stage-B fallback sensitivity.}
\textit{Purpose:} examine the behaviour of the Stage-B fallback that switches from the
flagged epochs to the full valid-epoch set when too few epochs are flagged.
\textit{Motivation:} the fallback count \(n_{\min}\) guards against estimating a threshold
from an insufficient sub-sample, but the choice of \(n_{\min}\) should not unduly drive the
results.
\textit{Objective:} quantify how threshold variance depends on the sub-sample size and how
often the fallback is triggered as a function of epoch duration, characterising the
estimator's stability under small flagged sets.

\paragraph{Experiment~6 --- Blink-region morphology (butterfly).}
\textit{Purpose:} visually and quantitatively characterise the morphology of the detected
blink regions.
\textit{Motivation:} aggregate $F_1$ does not reveal \emph{why} a detection succeeds or
fails; inspecting the waveforms distinguishes genuine, time-locked blink shapes from ragged
artefactual ones.
\textit{Objective:} produce butterfly overlays of true-positive, false-negative, and
false-positive blink regions, stratified by duration and amplitude, per subject and pooled,
to characterise the kinds of events the detector recovers and misses.

\paragraph{Experiment~7 --- Effect of excluding low-health epochs.}
\textit{Purpose:} measure how much the reported performance depends on the preprocessing
decision to exclude low-health epochs, for every detector.
\textit{Motivation:} dropping low-health epochs is a preprocessing step applied to all
methods, not part of any detector, so a reviewer will rightly ask whether the headline
numbers are partly an artefact of that exclusion.
\textit{Objective:} run every condition with and without the exclusion (health-on versus
health-off) on both datasets, reporting the per-condition change in $F_1$ to separate
preprocessing effects from detector effects and to test whether the proposed method depends
on the exclusion less than the baselines.

\paragraph{Experiment~8 --- Long-blink (drowsiness closure) recall.}
\textit{Purpose:} quantify recall specifically on long eye closures for every detector.
\textit{Motivation:} in a drowsy-driving context the long closures (microsleep,
\(\geq 0.5\)~s) are the safety-critical events, yet shape-based detectors tuned to short
blinks tend to reject wide-plateau closures.
\textit{Objective:} split blinks into normal (\(<0.5\)~s) and long (\(\geq 0.5\)~s)
categories and report the per-detector recall gap on both datasets, turning ``the method
also handles long closures'' from a claim into a measured result.


\subsection{Evaluation and Statistical Analysis}
Detection was scored at the event level: a detected interval that overlapped an unmatched ground-truth interval was counted as a true positive, unmatched detections were counted as false positives, and unmatched ground-truth intervals were counted as false negatives. Event-level precision, recall, and $F_1$ were macro-averaged across sessions, and event-level $F_1$ was preferred over accuracy because blink events were sparse relative to non-blink samples. Sensitivity to boundary tolerance was assessed across intersection-over-union thresholds {0.0, 0.1, 0.2, 0.3, 0.5}. Methods were compared using paired two-tailed Wilcoxon signed-rank tests on session-level $F_1$, with Bonferroni correction applied to six comparisons for the four-method pairwise tests and to four comparisons for the epoch sweep against the 30 s reference, for which the corrected alpha was 0.0125; rank-biserial effect sizes were also reported.

